\scanpage{13}
\deleted{Considérons un point $\xi\in S(k)$, \inserted{$k$ corps de car. $0$,} il définit un foncteur fibre $H^{**}_\xi$ (Hodge) et $H^*_\xi$ (De Rham) à valeurs dans la catégorie des vectoriels sur $k$.}

\begin{center}\rule{3cm}{0.4pt}\end{center}

Cas d'un \uline{tore} $T$ sur $\mathbb Q$.

Défini par un $\pi$-module admissible, libre de t.f. sur $\mathbb Z$, soit $M$, $(\pi=\pi_1(\mathbb Q,C))$

On aura $\varepsilon\in M^\pi$, $i\deleted{\unsure{\in}}\in\check M^\pi$, $\langle i,\varepsilon\rangle=2$.

Les modules sur $T$ correspondent aux éléments invariants par $\pi$ de $\mathbb Z(M)^+$ (comb. lin. à coeff. $\geq 0$)

\deleted{Donnons nous \unsure{le|$\varepsilon$} un}

La donnée d'une structure \deleted{\unsure{positive}} \inserted{de polarisation} ne dépend que du tore réel $T_{\mathbb R}$. Ce dernier est de la forme \deleted{\unsure{$T_2$|$T_{\mathbb R}\simeq$} $U\times S$,} \inserted{ext. d'un $S$ \unsure{splitté|split} par} $U$ compact, \inserted{i.e. $U(\mathbb R)$ cpct,} \deleted{$S$ \unsure{splitté|split},} et $\varepsilon$ est néc. triviale sur $U$. Donc on en déduit des structures $\varepsilon_S$ et $j_S$ sur $S$, et évidemment $S$ \uline{hérite} de structure de positivité. J'affirme qu'on doit avoir $\dim S=1$, en d'autres termes, que \uline{$\operatorname{Ker}\varepsilon$ est compact}.

En effet, \deleted{si on avait $\dim S\geq 2$, on \deleted{pourrait} peut \deleted{trouver un} se ramener (divisant par \deleted{\unsure{qq}} un sous-tore de codim 1 de $\operatorname{Ker}\varepsilon$) que la dim est deux. \deleted{Soit} On \deleted{\unsure{aura}} peut écrire $S=G_{m\,\mathbb R}\times G_{m\,\mathbb R}$, avec $i(\deleted{\unsure{t|1}}\lambda)=(\unsure{1|0|\lambda^\beta},\lambda^\alpha)$, $\varepsilon(\lambda,\mu)=\mu^{\unsure{\beta}}$, où $\alpha,\beta$ entiers $\geq 1$, $\alpha\beta=2$ (donc $\alpha\unsure{=|,\ }1$, $\beta=2$, ou $\alpha=2$, $\beta=1$).}
% archive sheet number: 12
